\chapter{Resultados Esperados}
A seção de Resultados Esperados visa apresentar as projeções e contribuições que a pesquisa pretende alcançar. Estes resultados não representam as conclusões finais, mas são estimativas baseadas nos objetivos definidos. Nesta seção, os resultados esperados são descritos de forma clara, objetiva e fundamentada sob a ótica da infraestrutura de informação.

\section{Introdução aos Resultados Esperados}
Os resultados esperados desta pesquisa estão diretamente relacionados ao desenvolvimento de uma infraestrutura orquestrada para laboratórios de Física. Pretende-se que os resultados contribuam para o avanço do conhecimento na área de Sistemas de Informação Aplicados à Educação e ofereçam uma solução tecnológica para o problema do isolamento e ociosidade das bancadas experimentais. O sucesso do artefato será medido pela sua capacidade de centralizar telemetria real em uma interface web estável.

\section{Descrição dos Resultados Esperados}
Os resultados esperados são organizados em dimensões técnicas (quantitativas) e estruturais (qualitativas), focando na viabilidade da arquitetura de Sistemas Embarcados proposta.

\subsection{Resultados Quantitativos}
Os resultados quantitativos incluem métricas de desempenho sistêmico e fidelidade de dados, esperadas como desdobramento da implementação da Prova de Conceito (PoC):
\begin{enumerate}
	\item Sincronização de Um Nó Escalável: Espera-se que a arquitetura suporte o monitoramento de uma bancada experimental (nó ESP32) em uma arquitetura descentralizada.
	\item Resolução de Telemetria: Obtenção de curvas de carga e descarga de capacitor, permitindo a identificação de leituras.
\end{enumerate}

\subsection{Resultados Qualitativos}
Os resultados qualitativos referem-se ao avanço no gerenciamento do ambiente laboratorial, proporcionando:
\begin{enumerate}
	\item Orquestração Centralizada: A partir da demonstração do funcionamento, mostrar a possibiliade da não necessidade de computadores individuais por bancada experimental em um laboratório de Física, com centralização da gestão dos dados.
	\item Persistência e Rastreabilidade: Criação de um histórico digital dos experimentos em banco de dados.
	\item Desacoplamento Tecnológico: Demonstração de uma arquitetura onde a interface do usuário é independente do hardware de coleta, facilitando futuras manutenções e inclusão de novos tipos de experimentos físicos.
\end{enumerate}

\section{Impactos Práticos e Científicos}
Os resultados esperados possuem implicações profundas na modernização da infraestrutura acadêmica:
\begin{enumerate}
	\item Impactos Práticos: Espera-se que a pesquisa forneça um modelo de laboratório escalável para instituições públicas, onde a comunicação sem fio (Wi-Fi) reduz o emaranhado de cabos e a complexidade de montagem física das experiências.
	\item Impactos Científicos: com a criação de ambientes de ensino com experimentação sem fio e descentralizada, espera-se que a arquitetura ofereça uma infraestrutura robusta que serve de alicerce para possíveis futuros Laboratórios de Acesso Remoto (LAR).
\end{enumerate}

\section{Fundamentação dos Resultados}
Os resultados esperados são fundamentados na robustez dos estudos prévios de arquiteturas web para ensino de Física associado a instrumentação experimental. Autores como Kuriščák et al. (2022) indicam que o desacoplamento entre UI e hardware é a chave para a escalabilidade de laboratórios integrados. Estudos anteriores de Vidal et al. (2022) reforçam que o uso de microcontroladores para coleta de dados reais — em oposição a simulações — aumenta o engajamento e a percepção de relevância científica pelos alunos. A viabilidade técnica de superar especificações comerciais com hardware de baixo custo é suportada pelos achados de Nino et al. (2014).

\section{Conclusão}
Diante do exposto, os resultados esperados para esta pesquisa objetivam a apresentação de uma infraestrutura técnica que suporte a orquestração de dados em laboratórios de Física. Espera-se que o artefato desenvolvido demonstre a viabilidade de centralizar a telemetria experimental de forma sem fio, podendo reduzir a complexidade logística das bancadas e permitindo um monitoramento mais integrado. Assim, a arquitetura proposta apresenta-se como uma alternativa viável para a modernização da gestão laboratorial.
